\documentclass[12pt,pdftex]{article}
\usepackage{amsmath}
\usepackage{graphicx}
\usepackage{natbib}
\citestyle{elsart-harv}
\usepackage{rotating}
\usepackage{lineno} 
\usepackage{longtable}
\usepackage{amsmath}
\usepackage{amssymb}


% Text layout
\topmargin 0.0cm
\oddsidemargin 0.5cm
\evensidemargin 0.5cm
\textwidth 16cm 
\textheight 21cm

% Leave date blank
\date{}

\title{Response to referee comments\\The Dimensionality of Ecological Networks}
                                                                  
\begin{document}
\linenumbers
\maketitle
\baselineskip 18pt

\noindent Dear Prof. Franklin,
THANK YOU FOR 

\noindent {\bf Comments to Author:}\\
This ms contributes to the effort to advance species distribution modeling approaches by considering the interplay of the velocity of climate change, dispersal ability of species, configuration of available area for the species, and demographic aspects. The ms finds that the species dispersal ability had the largest effect on both the fraction of population persisting after the climatic shift, and the global population density. Nonetheless, in terms of the fraction of occupied habitat patches, the distribution of habitat was also important. The authors conduct a theoretical analysis but also test their approach against 
four real world species in Europe. \\

The ms was reviewed by a single reviewer, and due to unavailability of several reviewers approached, I myself have provided in-depth comments below.\\

\noindent RESPONSE REF. \#1:\\
We thank the editor for the work he put into our manuscript ensuring it got rapidly handled. 

It is particularly useful to explore how climate change interacts with human transformation of natural habitat given that the latter is ubiquitous in most parts of the world, and growing in those parts of the world where this is less important. This topic has been addressed over more than a decade by a number of researchers, and I agree with the reviewer that not enough attention is given to earlier work. For example, Collingham and Huntley worked on this challenge more than a decade ago (MIGRATE) and there are many other examples. As a consequence, the ms feels as though the authors are reinventing the wheel to some extent. For example, it would be interesting to see discussion by the authors of how the utility of this approach compares with that of dispersal chains developed by Williams et al in 2005 and advanced by Phillips et al in several papers.\\

The approach to Landscape Generation is interesting and useful, and has important practical ramifications for conservation approaches. I support the reviewer suggestion that reference to the SLOSS debate would be valuable context, especially in discussion and conclusions.\\

The approach to dispersal kernel is also useful, though once again not referencing earlier work (e.g, Nathan, and Higgins et al), in which much effort was made to consider random long distance dispersal events and fat –tailed kernels. This aspect needs addressing, at the very least to acknowledge the issue, and to reflect on findings relative to this earlier work that suggested long distance dispersal events can vastly increase the ability of a species to track climate change, particularly across fragmented landscapes. Overall, it would be valuable for the authors to identify threshold values of landscape configuration that match different dispersal abilities, to start o provide quantitative guidance to conservation planners on the kind of effort needed to provide suitable habitat.\\

I also agree with the reviewer that the initial inclusion of 8 species, followed by the abandonment of four of them because their ranges contracted with climate change, is confusing. Why not just leave them out? \\

This major suggestion by the reviewer bears consideration, to quote “Could you not just make this paper about the interaction of landscape, dispersal distance (and demographics) and climate velocity? If these are theoretical constructs, then you don’t need such a long methods section, and you can still relate the results to potential target species in the discussion. If on the other hand, you decide this is more about the species that were modeled, and there are specific conservation messages about those species (as opposed to generalized results that also have indications), then you might consider using HexSim, and inputting actual climate values for current and projected futures. The projected futures are not likely to show a similar “moving window” of climate velocity, but to exhibit more varied change on the landscape. I would suggest going with the former, more theoretical approach.”\\

My sense of the suggestion above is that if you simply leave out the section covering actual species modeled, the paper as is loses a lot of novelty and retreads ground previously explored. A more thorough referencing of the literature might be useful for the authors to demonstrate this novelty more strongly, and inform better the choice of how to deal with melding the two main thrusts of this paper together, or indeed focusing more on the theoretical aspects. \\

It is my recommendation that the authors consider the suggestions and criticisms above, together with the comments of the reviewer, and revise the ms accordingly, especially taking into account the literature in this field over the past decade, and especially in the latter period, including a paper by Thorne et al, published in this journal recently, along similar lines.\\

Professor Guy Midgley



\begin{flushright}
  Best Regards,\\
  Anna Ekl\"of
\end{flushright}

\subsection*{Reviewer 1 -- Dr. Gravel}

\end{document}
